\chapter{Development and extension}
\label{ch:development}

This chapter is for people extending the program: adding a rule, a parser
section or an analyser, or packaging a release.

\section{Architecture and the layering contract}

\begin{codeblock}
ui/            menu, command line, script replay, session
  recipe/      system profile -> recommendations -> input templates
  diagnosis/   rule engine (findings with severity and provenance)
  analysis/    characterisers (output files, structures, data files)
  toolindex/   the external-tool index
      parsers/     program-output parsing (ORCA deep)
        knowledge/ rules, basis/ECP data, templates, sources.bib
\end{codeblock}

Dependencies point one way only, and the direction is enforced mechanically by
import-linter (the first step of \file{scripts/verify.sh}): ui may use
everything below it; recipe/diagnosis/analysis/toolindex may use parsers and
knowledge; parsers may use knowledge; knowledge uses nothing from the program.
Two house rules accompany the layering: diagnosis and recipe only ever
\emph{propose} (they write new files, never modify yours), and analysis never
depends on diagnosis (it must work on any output file, stand-alone).

\section{Adding a rule}

\begin{enumerate}
  \item Register any new fact field in \file{knowledge/rules/README.md} first,
    and make the parser layer produce it (\file{parsers/orca.py}, function
    \code{facts()}); add an evaluation case to \file{tests/test\_knowledge.py}.
  \item Add the rule to a YAML file under \file{knowledge/rules/}. Every rule
    needs non-empty evidence; a \code{literature} item must point into
    \file{sources.bib} with a bibkey (the loader rejects it otherwise); a
    \code{diagnosis} rule needs a severity and no \code{method}; a
    \code{recipe} rule needs a \code{method} and no severity.
  \item Quote evidence text as a YAML block scalar (\code{>-}). A
    double-quoted scalar ends at the closing quote, and nothing may follow on
    the line -- the classic trap this project's rule files document.
  \item Add tests: the rule logic against hand-built fact tables, plus an
    end-to-end case on a real fixture.
\end{enumerate}

\section{Adding a parser section}

The fixture discipline comes first: obtain a real output file that exercises
the pattern, register it in \file{fixtures/orca/README.md} with its source and
run outcome, and only then write the regular expressions against it. Manual
examples may come from older program versions and are not authoritative. New
patterns must not fire on the existing fixtures incorrectly: run the whole
suite, and check negative cases explicitly (a marker that is always present is
unusable -- the SOC detection in this project was fixed by falsifying candidate
markers one by one against five fixtures).

\section{Adding an analyser}

An output-reading analyser implements \code{accepts(result)} and
\code{run(result) -> ReportSection} and provides \code{evidence()}; it is
registered in \file{analysis/\_\_init\_\_.py}, after which the menu-1 pipeline
runs it automatically and merges its evidence into the report's provenance.
Structure-reading and mapping-reading analysers (geometry, point charge,
declaration check) are called directly by their menu handlers instead.
Everything user-visible is in English; error messages end with a
``Next step: \dots'' sentence.

\section{Verification and packaging}

\begin{codeblock}
bash scripts/verify.sh
\end{codeblock}

runs, in order: the layering contract; a byte-compile of all sources and tests
with the interpreter at hand (this is the syntax gate for the supported
minimum, Python 3.11 -- development machines often run newer interpreters that
accept syntax 3.11 rejects); and the test suite. The suite includes the
examples of this manual: \file{tests/test\_manual.py} replays every script in
\file{docs/manual/examples/} and compares the output byte for byte.

Packaging produces two artefacts (Chapter~\ref{ch:installation}): the pip wheel
(\code{pip wheel .}) and the green package (PyInstaller with
\file{packaging/fblockkit.spec}, built inside a 3.12 environment). Two
packaging lessons are baked into the spec and worth repeating: every data-file
suffix must be declared both in \file{pyproject.toml} (package-data) and in the
spec's collection pattern (a missing \file{.bib} once broke the green package
while the wheel worked); and always \emph{run} the packaged build, not just
build it.

\section{The examples of this manual}

\file{docs/manual/examples/} contains the session scripts, the deterministic
input generator and the captured outputs; \file{run-examples.sh} refreshes the
captures. When a change alters any program output, refresh the captures and
re-run the suite: a stale capture is a failing test, by design.
